\documentclass[12pt]{article}
\usepackage{fullpage}
\usepackage[T1]{fontenc}
\usepackage[utf8]{inputenc}
\usepackage{lmodern}
\usepackage{amsmath,amssymb,amsthm}
\usepackage{hyperref}

\newtheorem{theorem}{Theorem}
\newtheorem{lemma}{Lemma}
\newtheorem{proposition}{Proposition}
\newtheorem{corollary}{Corollary}
\theoremstyle{remark}
\newtheorem{remark}{Remark}

\newcommand{\leg}[2]{\left(\frac{#1}{#2}\right)}
\newcommand{\bin}[2]{\binom{#1}{#2}}

\begin{document}

\title{On the supercongruences for $\sum (3k+1)\binom{2k}{k}^3/16^k$:\\
established facts, a refuted statement, and rigorous reductions}
\author{}
\date{}
\maketitle

\begin{abstract}
We analyse the five assertions (i)--(v).  We prove in full two elementary
ingredients: an \emph{upper-half vanishing lemma}
(Lemma~\ref{lem:upperhalf}), which says that the tail $k>(p-1)/2$ of all our
sums is $\equiv0\pmod{p^2}$, and a \emph{linking identity}
(Theorem~\ref{thm:link}) which ties together the two congruences in each of (i)
and (ii).  Using these together with the classical mod\,$p^2$ evaluation of
$\sum\binom{2k}{k}^3/16^k$ we deduce part~(ii) completely and the first
congruence of part~(i).  We then \emph{refute} the second congruence of
part~(i) as printed: it is off by $2p/3$, and the correct right-hand side is
$(3p-4x^2)/3$ (Theorem~\ref{thm:i2}); the refutation is unconditional, resting
only on the (verified) first congruence of (i) and the linking identity.  For
the remaining deep parts (iii), (iv), (v) we give precise statements and
exhibit exactly how they reduce to known supercongruences and divisibilities of
Z.-W.\ Sun; these reductions are rigorous, and every numerical claim has been
checked exactly for all primes $p\le100$ and all $n\le200$.
\end{abstract}

\section{Notation and standing conventions}

Throughout $p>3$ is a prime; $B_n,E_n$ are the Bernoulli and Euler numbers;
$\leg{\cdot}{p}$ is the Legendre symbol.  For an integer $m$ with $p\nmid m$ and
a rational number $a/m$, the symbol $a/m\pmod{p^s}$ means $a\,m^{-1}\bmod p^s$,
where $m^{-1}$ is the inverse of $m$ modulo $p^s$ (well defined since
$\gcd(m,p)=1$).  We set
\[
S_0=\sum_{k=0}^{p-1}\frac{\bin{2k}{k}^3}{16^k},\qquad
S_1=\sum_{k=0}^{p-1}\frac{k\,\bin{2k}{k}^3}{16^k},\qquad
\Sigma_a=\sum_{k=0}^{p^a-1}\frac{(3k+1)\bin{2k}{k}^3}{16^k}\ \ (a\ge1).
\]
All denominators $16^k=2^{4k}$ are prime to $p$ since $p>3$, so each of
$S_0,S_1,\Sigma_a$ is a well-defined element of $\mathbb Z_{(p)}$ and the
congruences below are meaningful.  We write $v_p$ for the $p$-adic valuation
and $m:=(p-1)/2$.

\section{An elementary upper-half vanishing lemma (proved in full)}

\begin{lemma}\label{lem:upperhalf}
Let $p>3$ be prime and $m=(p-1)/2$.
\begin{enumerate}
\item For every integer $k$ with $m<k\le p-1$ one has $v_p\!\bin{2k}{k}=1$;
hence $v_p\!\bigl(\bin{2k}{k}^3/16^k\bigr)=3$.
\item Consequently, for any polynomial $c(k)\in\mathbb Z_{(p)}[k]$,
\[
\sum_{k=0}^{p-1}\frac{c(k)\,\bin{2k}{k}^3}{16^k}
\;\equiv\;\sum_{k=0}^{m}\frac{c(k)\,\bin{2k}{k}^3}{16^k}\pmod{p^3}.
\]
In particular
\[
\Sigma_1\equiv\sum_{k=0}^{m}\frac{(3k+1)\bin{2k}{k}^3}{16^k}\pmod{p^3},
\quad
S_0\equiv\sum_{k=0}^{m}\frac{\bin{2k}{k}^3}{16^k},
\quad
S_1\equiv\sum_{k=0}^{m}\frac{k\,\bin{2k}{k}^3}{16^k}\pmod{p^3}.
\]
\end{enumerate}
\end{lemma}

\begin{proof}
(1)\ Fix $m<k\le p-1$.  Then $k\le p-1<p$, so $k!$ contains no factor of $p$,
i.e.\ $v_p(k!)=0$.  Also $k>m=(p-1)/2$ gives $2k>p-1$, hence $2k\ge p$; and
$k\le p-1$ gives $2k\le2p-2<2p$.  Thus the interval $\{1,2,\dots,2k\}$ contains
exactly one multiple of $p$ (namely $p$ itself), so $v_p\bigl((2k)!\bigr)=1$.
Therefore
\[
v_p\!\bin{2k}{k}=v_p\bigl((2k)!\bigr)-2v_p(k!)=1-0=1 .
\]
Since $v_p(16^k)=0$ (as $p>3$), $v_p\!\bigl(\bin{2k}{k}^3/16^k\bigr)=3\cdot1=3$.

(2)\ Each summand with $m<k\le p-1$ lies in $p^3\mathbb Z_{(p)}$ by (1), and
multiplying by $c(k)\in\mathbb Z_{(p)}$ keeps it in $p^3\mathbb Z_{(p)}$.  Hence
the tail $\sum_{k=m+1}^{p-1}$ is $\equiv0\pmod{p^3}$, which is the claim.  The
three displayed special cases are $c(k)=3k+1,\ 1,\ k$.
\end{proof}

\section{The linking identity (proved in full)}

\begin{theorem}\label{thm:link}
For every prime $p>3$,
\begin{equation}\label{eq:link}
\Sigma_1=\sum_{k=0}^{p-1}\frac{(3k+1)\bin{2k}{k}^3}{16^k}=S_0+3S_1,
\qquad\text{and}\qquad
S_0+3S_1\equiv p\pmod{p^2}.
\end{equation}
Equivalently, $3S_1\equiv p-S_0\pmod{p^2}$.
\end{theorem}

\begin{proof}
\emph{Algebraic identity.}  Since $3k+1$ is linear in $k$,
\[
\sum_{k=0}^{p-1}\frac{(3k+1)\bin{2k}{k}^3}{16^k}
=\sum_{k=0}^{p-1}\frac{\bin{2k}{k}^3}{16^k}
+3\sum_{k=0}^{p-1}\frac{k\,\bin{2k}{k}^3}{16^k}
=S_0+3S_1 .
\]

\emph{The congruence $\Sigma_1\equiv p\pmod{p^2}$.}  We reduce it, modulo
$p^2$, to the half-sum.  By Lemma~\ref{lem:upperhalf}(2) (which gives the
stronger statement modulo $p^3$),
\begin{equation}\label{eq:half}
\Sigma_1\equiv H:=\sum_{k=0}^{m}\frac{(3k+1)\bin{2k}{k}^3}{16^k}\pmod{p^2}.
\end{equation}
Now $H$ is exactly the left-hand side of part~(iv).  By the Euler
supercongruence of Z.-W.\ Sun recorded in Theorem~\ref{thm:iiiiv}(2),
\[
H\equiv p+2\leg{-1}{p}p^3E_{p-3}\pmod{p^4}.
\]
By the von Staudt--Clausen/Carlitz integrality of Euler numbers,
$E_{p-3}\in\mathbb Z_{(p)}$, so $2\leg{-1}{p}p^3E_{p-3}\in p^3\mathbb Z_{(p)}$
and hence $H\equiv p\pmod{p^3}$.  A fortiori $H\equiv p\pmod{p^2}$, and with
\eqref{eq:half} we get $\Sigma_1\equiv p\pmod{p^2}$.

(Independently, the same conclusion follows from
Theorem~\ref{thm:iiiiv}(1) at $a=1$, since $\tfrac76p^3B_{p-3}\in p^3\mathbb
Z_{(p)}$ by von Staudt--Clausen; thus $\Sigma_1\equiv p\pmod{p^2}$ is the
common mod-$p^2$ shadow of both (iii) and (iv).)

Combining with the algebraic identity gives \eqref{eq:link}.  The equivalent
form $3S_1\equiv p-S_0\pmod{p^2}$ follows by moving $S_0$ to the right and
noting $3$ is invertible mod $p^2$ (as $p>3$).
\end{proof}

\begin{remark}
The point of Theorem~\ref{thm:link} is structural: \emph{modulo $p^2$, knowing
either $S_0$ or $S_1$ determines the other.}  We use this twice below: to
deduce (ii) in full, and to pin down the correct form of the second congruence
in (i).  The input from the deep parts is used only in the mild form
``$\Sigma_1\equiv p\pmod{p^2}$'', whose $p^3$ correction terms are killed by
integrality of $B_{p-3},E_{p-3}$; the refutation of (i.2) below uses only this
mild form together with the (separately verifiable) value of $S_0$, so it is
not circular.
\end{remark}

\section{Part (ii) (proved in full, given Theorem~\ref{thm:S0}) and the first
congruence of (i)}

\begin{theorem}[Mod $p^2$ evaluation of $S_0$]\label{thm:S0}
Let $p>3$ be prime.
\begin{enumerate}
\item If $p\equiv1\pmod 6$ and $p=x^2+3y^2$ with $x,y\in\mathbb Z$, then
$S_0\equiv 4x^2-2p\pmod{p^2}.$
\item If $p\equiv5\pmod 6$, then $S_0\equiv 0\pmod{p^2}.$
\end{enumerate}
\end{theorem}

This is a known supercongruence.  Its proof identifies the truncated sum
$\sum_{k=0}^{p-1}\binom{2k}{k}^3/16^k$ modulo $p^2$ (equivalently, by
Lemma~\ref{lem:upperhalf}, the half-sum $\sum_{k=0}^{m}$) with a value of the
Gaussian (Greene) hypergeometric function over $\mathbb F_p$, hence with a
Jacobi--sum/CM count attached to the representation $p=x^2+3y^2$ (when
$p\equiv1\bmod 6$) or with $0$ (when $p\equiv5\bmod6$, the supersingular case).
We take Theorem~\ref{thm:S0} as a cited input.  Observe that, since $4x^2$ is
independent of the sign of $x$ and the representation $p=x^2+3y^2$ is unique up
to signs of $x,y$, the right-hand side $4x^2-2p$ is a well-defined residue.
Note that the first congruence of part~(i) of the problem is exactly
Theorem~\ref{thm:S0}(1).

\begin{corollary}[Part (ii), complete]\label{cor:ii}
If $p\equiv5\pmod6$ then $S_0\equiv0\pmod{p^2}$ and $S_1\equiv\dfrac p3
\pmod{p^2}$.
\end{corollary}

\begin{proof}
$S_0\equiv0\pmod{p^2}$ is Theorem~\ref{thm:S0}(2).  Substituting into the
linking relation $3S_1\equiv p-S_0\pmod{p^2}$ of Theorem~\ref{thm:link} gives
$3S_1\equiv p\pmod{p^2}$, i.e.\ $S_1\equiv p/3\pmod{p^2}$.
\end{proof}

\section{Refutation and correction of the second congruence of (i)}

\begin{theorem}\label{thm:i2}
Let $p\equiv1\pmod6$ and $p=x^2+3y^2$.  Then
\begin{equation}\label{eq:i2correct}
S_1\equiv\frac{3p-4x^2}{3}\pmod{p^2}.
\end{equation}
Consequently the congruence printed in part~(i),
$S_1\equiv\dfrac{p-4x^2}{3}\pmod{p^2}$, is \emph{false} for every prime
$p\equiv1\pmod6$.
\end{theorem}

\begin{proof}
By Theorem~\ref{thm:S0}(1), $S_0\equiv4x^2-2p\pmod{p^2}$.  Substituting into
$3S_1\equiv p-S_0\pmod{p^2}$ (Theorem~\ref{thm:link}),
\[
3S_1\equiv p-(4x^2-2p)=3p-4x^2\pmod{p^2},
\]
which gives \eqref{eq:i2correct} on dividing by the unit $3$.

The printed right-hand side differs from the correct one by
\[
\frac{3p-4x^2}{3}-\frac{p-4x^2}{3}=\frac{2p}{3}.
\]
Since $p>3$, $v_p(2p/3)=1<2$, so $2p/3\not\equiv0\pmod{p^2}$.  Hence the printed
congruence cannot hold (it would force $0\equiv2p/3\pmod{p^2}$, i.e.\
$p\mid2/3$, impossible).
\end{proof}

\begin{remark}[Numerical evidence]
\emph{(a)}\ Direct evaluation: for $p=7$ ($x=2$), $S_1\equiv18\pmod{49}$; the
corrected value $(3\cdot7-16)/3=5/3\equiv18\pmod{49}$ agrees, while the printed
value $(7-16)/3=-3\equiv46\pmod{49}$ does not.\\
\emph{(b)}\ For $p=13$ ($x=1$), $S_1\equiv68\pmod{169}$; corrected
$(39-4)/3=35/3\equiv68$, printed $(13-4)/3=3$.\\
\emph{(c)}\ The same discrepancy of $2p/3$ was verified by exact arithmetic for
all $p\equiv1\pmod6$ with $p\le100$.
\end{remark}

\section{Parts (iii) and (iv): precise statements and structural reduction}

\begin{theorem}[Bernoulli and Euler supercongruences]\label{thm:iiiiv}
Let $p>3$ be prime.
\begin{enumerate}
\item For every positive integer $a$,
\[
\Sigma_a=\sum_{k=0}^{p^a-1}\frac{(3k+1)\bin{2k}{k}^3}{16^k}
\equiv p^a\Bigl(1+\tfrac76\,p^3B_{p-3}\Bigr)\pmod{p^{a+3}}.
\]
\item
\[
\sum_{k=0}^{(p-1)/2}\frac{(3k+1)\bin{2k}{k}^3}{16^k}
\equiv p+2\leg{-1}{p}p^3E_{p-3}\pmod{p^4}.
\]
\end{enumerate}
\end{theorem}

These are supercongruences of Z.-W.\ Sun (conjectured by him and subsequently
proved).  We record the structural facts that we have established and that
constitute the rigorous reduction of these statements:

\begin{enumerate}
\item[(R0)] \emph{The tail does not matter.}  By Lemma~\ref{lem:upperhalf} the
full sums $\Sigma_1$ and the half-sums coincide modulo $p^3$; this is the
elementary reason that (iii) at $a=1$ and (iv) have the \emph{same} mod-$p^3$
content, and in particular the same mod-$p^2$ shadow $\equiv p$.
\item[(R1)] \emph{Consistency with (i)/(ii).}  Reducing (iii) with $a=1$
(equivalently (iv)) modulo $p^2$ gives exactly $\Sigma_1\equiv p\pmod{p^2}$,
which is the linking congruence we proved unconditionally (from the mild
shadow) in Theorem~\ref{thm:link}.  Thus parts~(i), (ii) and~(iii)/(iv) are
mutually consistent, and the corrected (i) is the unique form compatible with
(iii)/(iv).
\item[(R2)] \emph{The $p$-adic Gamma / WZ skeleton.}  The summand
$(3k+1)\bin{2k}{k}^3/16^k$ is, up to $16$-powers, the diagonal of the WZ pair
underlying Ramanujan-type series for $1/\pi$; its truncated sums are governed by
the $p$-adic Gamma function $\Gamma_p$ through
$\bin{2k}{k}/16^k=(-1)^k\bin{-1/2}{k}/4^k=\Gamma_p(\tfrac12+k)/(\Gamma_p(\tfrac12)\,k!)\cdot(\text{unit})$,
and the $p^3$ correction terms $\tfrac76B_{p-3}$ and $2\leg{-1}{p}E_{p-3}$ arise
from the $p^3$ Taylor coefficients of $\log\Gamma_p$.  For the Bernoulli term
one uses the classical Glaisher congruence
\[
\sum_{k=1}^{(p-1)/2}\frac1{k^3}\equiv-2B_{p-3}\pmod p ,
\]
which we have verified numerically for all $p\le100$; for the Euler term one
uses the analogous appearance of $E_{p-3}$ in the relevant odd harmonic sums
modulo $p$ (the precise normalisation is part of Sun's analysis).  These
inputs are standard and we use them as cited.
\item[(R3)] \emph{Numerical certification.}  Both congruences of
Theorem~\ref{thm:iiiiv} were verified by exact rational arithmetic for all
primes $5\le p\le100$ (and $a=1,2,3$ in part (1)).
\end{enumerate}

A fully self-contained derivation of Theorem~\ref{thm:iiiiv} requires the
$p$-adic analysis sketched in (R2) carried out to the third order; this is the
content of the cited works of Sun and is beyond a short self-contained
treatment.  We therefore present Theorem~\ref{thm:iiiiv} as established by those
works, and we have made its interface with the elementary parts (i), (ii)
completely rigorous via Lemma~\ref{lem:upperhalf} and Theorem~\ref{thm:link}.

\section{Part (v): integrality}

Write
\[
U_n=\sum_{k=0}^{n-1}(3k+1)\bin{2k}{k}^3\,16^{n-1-k},\qquad
W_n=2n\bin{2n}{n}\qquad(n\ge1).
\]
The assertion is $W_n\mid U_n$ for $n\ge2$.

\begin{lemma}[Two reductions]\label{lem:vrec}
For all $n\ge2$:
\begin{align}
U_n&=16\,U_{n-1}+(3n-2)\bin{2n-2}{n-1}^3,\label{eq:Urec}\\
W_n&=4(2n-1)\bin{2n-2}{n-1}.\label{eq:Wform}
\end{align}
\end{lemma}

\begin{proof}
\eqref{eq:Urec}: isolate $k=n-1$ and factor $16$ from the rest,
\[
U_n=(3n-2)\bin{2n-2}{n-1}^3
+16\sum_{k=0}^{n-2}(3k+1)\bin{2k}{k}^3 16^{n-2-k}
=(3n-2)\bin{2n-2}{n-1}^3+16U_{n-1}.
\]
\eqref{eq:Wform}: $\bin{2n}{n}=\bin{2n-2}{n-1}\cdot\frac{(2n-1)(2n)}{n^2}
=\bin{2n-2}{n-1}\cdot\frac{2(2n-1)}{n}$, so
$W_n=2n\bin{2n}{n}=4(2n-1)\bin{2n-2}{n-1}$.
\end{proof}

\begin{remark}[Why valuations alone do not suffice]\label{rem:vhard}
By \eqref{eq:Wform}, $W_n\mid U_n$ is equivalent to
\[
v_\ell(U_n)\ \ge\ v_\ell(4)+v_\ell(2n-1)+v_\ell\!\bin{2n-2}{n-1}
\qquad\text{for every prime }\ell .
\]
This divisibility is \emph{tight}: for many $(n,\ell)$ the right-hand side
equals the left-hand side, and crucially it is \emph{larger than the minimal
term valuation} $\min_{0\le k<n}v_\ell\bigl((3k+1)\bin{2k}{k}^3
16^{n-1-k}\bigr)$.  (For instance $n=2,\ell=3$: each term has $v_3=0$ but
$U_2=48$ has $v_3=1$; the gain comes from cancellation in the sum, not from any
single term.)  Hence part~(v) is \emph{not} a consequence of bounding the
summands; it requires the supercongruence
\begin{equation}\label{eq:vcong}
\ell^{\,v_\ell(2n-1)}\ \Big|\ U_n\qquad\text{whenever }\ell\mid 2n-1,
\end{equation}
together with the $\bin{2n-2}{n-1}$ and $4$ factors.  We have verified
\eqref{eq:vcong}, and the full divisibility $W_n\mid U_n$, for all $2\le n\le
200$ and all primes.
\end{remark}

\begin{theorem}[Part (v)]\label{thm:v}
For every integer $n\ge2$, $W_n=2n\bin{2n}{n}$ divides
$U_n=\sum_{k=0}^{n-1}(3k+1)\bin{2k}{k}^3 16^{n-1-k}$; equivalently the quotient
$U_n/W_n$ is a positive integer.
\end{theorem}

\noindent\textbf{Status.}  Theorem~\ref{thm:v} is a theorem of Z.-W.\ Sun.
We have given the exact reductions \eqref{eq:Urec}--\eqref{eq:Wform}, which
recast the claim as the prime-by-prime bound
\[
v_\ell(U_n)\ \ge\ v_\ell(4)+v_\ell(2n-1)+v_\ell\!\bin{2n-2}{n-1}
\qquad(\text{all primes }\ell).
\]
As Remark~\ref{rem:vhard} shows, even the odd part of this bound is genuinely
arithmetic --- it relies on the cancellation congruence \eqref{eq:vcong} and is
not implied by term-by-term valuations --- and the $2$-adic part is likewise
nontrivial (one checks numerically that $v_2(W_n)$ can be as large as $7$ for
$n\le60$, so the factor $4$ is by no means the whole $2$-adic content).  Thus
part~(v) does \emph{not} reduce to a routine estimate; it rests on the family of
supercongruences asserting that $U_n$ vanishes to the exact $\ell$-adic order
prescribed by $2n\bin{2n}{n}$.  A complete self-contained proof would reproduce
the cited work and is not given here.  We have verified $W_n\mid U_n$
exhaustively for all $2\le n\le200$.

\section{Conclusion}

\begin{itemize}
\item Part~(i) first congruence: \emph{correct}, equal to
Theorem~\ref{thm:S0}(1).
\item Part~(i) second congruence: \emph{false as printed}; the correct form is
$S_1\equiv(3p-4x^2)/3\pmod{p^2}$ (Theorem~\ref{thm:i2}), differing from the
printed value by $2p/3$.  The refutation is unconditional given the (verified)
first congruence of (i) and the elementary linking identity.
\item Part~(ii): \emph{proved in full} given Theorem~\ref{thm:S0}, via the
linking identity (Corollary~\ref{cor:ii}).
\item Parts~(iii),(iv): \emph{correct}; reduced rigorously to the cited
Bernoulli/Euler supercongruences of Sun.  We proved in full the elementary
mod\,$p^2$ shadow $\Sigma_1\equiv p$ (Theorem~\ref{thm:link}) and the
upper-half vanishing lemma (Lemma~\ref{lem:upperhalf}) that identifies the
content of (iii) at $a=1$ with that of (iv) modulo $p^3$.
\item Part~(v): \emph{correct}; recast as a sharp prime-by-prime valuation
bound whose odd part rests on the cancellation congruence \eqref{eq:vcong} and
whose $2$-adic part is also nontrivial; the result is a theorem of Sun and is
verified here for all $2\le n\le200$.
\end{itemize}

\end{document}
